\section{Related work}
\label{sec:related}

This report sits at the intersection of five literatures: the design of small LLM-based
\emph{guard classifiers}; the growing evidence that \emph{fine-tuning degrades} or narrows a
model's safety behavior; the study of \emph{calibration and operating points} for moderation;
\emph{model composition} in weight versus output space; and the emerging work
on \emph{domain-specialized guarding} in regulated verticals. We review each in turn and, for each,
name precisely what it establishes and what our paired, same-checkpoint,
composition-aware, four-domain design adds. Our contribution is not a new model, metric, or training
algorithm --- it is a measurement discipline and four evaluation instruments applied to one fixed
panel, so the contrast throughout is methodological rather than a claim of superior scores.

\begin{background}{How to read this map (for the non-expert)}
Almost all of the ``guard'' papers below ask \emph{how good is model $X$?} and answer with a single
benchmark number, usually next to a table of \emph{other} models. That is a \emph{leaderboard}
question. This report asks a different, paired question: \emph{what did the fine-tune change relative
to the same model before tuning, and does that change survive on data the guard never saw?} Keep that
distinction in mind --- most gaps we point to are gaps between a leaderboard number and a paired,
regime-split delta, not disagreements about which model is ``best.''
\end{background}

\subsection{Small LLM guard classifiers}
\label{sec:related-guards}
The dominant deployment pattern is a compact decoder LLM turned into a binary or taxonomy-tagged
moderation head. Llama Guard \citep{inan2023llamaguard} introduced the input--output
safeguard framing and a safety taxonomy, and later iterations pushed toward smaller, cheaper, and
multimodal variants \citep{meta2024llamaguard2,meta2024llamaguard3_8b,meta2024llamaguard3_1b,%
fedorov2024llamaguard3_1b,meta2025llamaguard4}. ShieldGemma \citep{zeng2024shieldgemma} and Granite
Guardian \citep{padhi2024graniteguardian} extend the recipe to other base families with broad harm
taxonomies; WildGuard \citep{han2024wildguard} adds one-stop coverage of prompt harm,
response harm, and refusal detection; and Qwen3Guard \citep{qwen3guard2025} is a recent open guard on
the same Qwen family two of our checkpoints come from. Beyond the general-purpose guards, several works
target narrower slices or richer inference: dedicated prompt-injection detectors
\citep{meta2024promptguard,meta2025promptguard2_86m,meta2025promptguard2_22m}, reasoning- and
logic-augmented guardrails \citep{kang2025r2guard}, ensemble-of-experts moderation
\citep{ghosh2024aegis}, RL-driven multilingual guardrails \citep{deng2025duoguard}, and CPU-class or
multi-stage pipelines aimed at cost \citep{majhi2025gpuguard}. The parameter-efficient adaptation we
use --- LoRA \citep{hu2021lora} applied to a chat model to produce a guard --- is exactly the LoRA-Guard
recipe \citep{elesedy2024loraguard}, and standardized suites such as GuardBench
\citep{bassani2024guardbench} have made cross-model comparison routine.

What this family reports, almost without exception, is \emph{absolute} moderation performance of
\emph{one} guard against \emph{different} models on a benchmark. That is precisely the quantity we argue
is co-produced by the benchmark and therefore uninformative about a specific fine-tune. Several of these
works do report in-distribution versus out-of-distribution blocks --- Llama~Guard's own test set against
zero-shot ToxicChat and OpenAI-Mod, WildGuard's and LoRA-Guard's ID/OOD splits --- so the
represented/transfer \emph{distinction} is not ours. Our narrower claim, stated once and not widened
anywhere else in this report, is:

\begin{quote}
a same-checkpoint \emph{paired} measurement of LoRA-induced change across represented and
source-held-out regimes, together with a calibrated base-plus-adapter composition test under matched
false-alarm budgets.
\end{quote}

\noindent An earlier version of this passage said instead that ``none of them reports'' the paired
change and that ``none treats retraining-free composition as a measurable design axis.'' Both were
overclaims that survived because the comparison was made in prose rather than against a table; the
per-component comparison in \Cref{tab:closest} is what replaces them. We reuse the LoRA-Guard recipe not to beat
these guards on a leaderboard --- we score four small open checkpoints, not the gated production guards
--- but to isolate what the fine-tune itself does. GuardBench and its peers are the backdrop against
which \Cref{sec:actI}'s question is posed: they normalize the single-suite comparison whose fragility we
show is fragile.

\subsection{Fine-tuning degradation and policy / benchmark transfer}
\label{sec:related-transfer}
The closest prior evidence is that fine-tuning can \emph{weaken} rather than strengthen safety.
\citet{hsiung2025guardrails} show that safety guardrails can collapse after fine-tuning and attribute
the collapse to similarity between the alignment and fine-tuning data --- a mechanism that rhymes with
our specialization finding, but is measured on a model's own alignment behavior rather than on a guard
classifier's represented-versus-transfer ranking split. \citet{liu2026domaingeneralizable} document that
guardrails overfit their training \emph{policy} and propose augmented policy training as a remedy;
\citet{li2026surfaceheuristics} give a complementary mechanistic account, showing that prompt-attack
defenses learn \emph{surface heuristics} that do not generalize --- essentially a why for the transfer
loss we observe empirically. \citet{bassani2025robustness} probe the same fragility from the
perturbation side, measuring guardrail robustness to input mutations and adversarial attacks, and
\citet{hackett2025bypassing} demonstrate concrete evasion of injection/jailbreak detectors. Framed most
generally, \citet{akinrele2026regime} argue that prompt-injection detection is \emph{regime-dependent} ---
performance is a function of the deployment regime, not a fixed model property --- which is our thesis
stated for one task with interpretable structural signals.

Our addition is the measurement design rather than the qualitative claim. Where these works compare a
model before and after, or one policy against another, we hold the \emph{checkpoint, manifest, seeds,
and scorer fixed} and read the paired change as a distribution over a purposively chosen panel with a
hierarchical bootstrap, and we decompose it into a large represented-source gain versus a near-flat but
\emph{heterogeneous} transfer change (\Cref{tab:primary}, \Cref{fig:plane}). Crucially, we do not stop
at ``degradation'': we quantify the specialization geometry per checkpoint and per benchmark, carry it
to a deployable operating point (\Cref{tab:sensitivity}), and then ask whether a composition
(\Cref{sec:actIII}) changes it. And unlike \citet{liu2026domaingeneralizable}, whose remedy retrains with
augmented policies, our candidate remedy retrains nothing.

\subsection{Calibration and operating points}
\label{sec:related-calibration}
Ranking quality (average precision) and thresholded decision quality are distinct, and the gap between
them is a calibration problem. \citet{guo2017calibration} established that modern neural networks are
systematically miscalibrated and that simple post-hoc scaling helps; \citet{liu2025calibration}
specialized this to LLM-based guard models, showing they are poorly calibrated for reliable content
moderation; and FlexGuard \citep{flexguard2026} moves past a single fixed threshold toward continuous,
strictness-adaptive risk scoring. This literature motivates two choices we make: we fit calibrators only
on a development split before reading any threshold, and we report operating-point behavior (macro- and
pooled-FPR, TPR, single-class recall) \emph{separately} from ranking.

Our specific contribution to this thread is a clean separation of the two failure modes on the \emph{same}
guards. In Act~II we show that a composition can recover transfer \emph{ranking} while its realized
false-alarm rate at a fixed FPR target still misses (\Cref{tab:pilot-operating}): recovering rank does
not deliver a transferable threshold. In Act~III the same lesson recurs from the other direction --- for
tightly clustered guard scores the fixed-threshold operating point is knife-edge and unstable across
software versions, which is why we deliberately do not tabulate it there. Where the calibration
literature asks ``is this score a probability?'', we use that machinery to make the sharper point that
\emph{ranking recovery $\neq$ calibration transfer}, a distinction a benchmark AP number alone hides.

\subsection{Weight-space versus output-space composition}
\label{sec:related-composition}
Combining models to improve robustness under distribution shift is well studied in \emph{weight} space:
WiSE-FT interpolates a fine-tuned model with its zero-shot initialization \citep{wortsman2022wiseft},
and model soups average the weights of several fine-tunes \citep{wortsman2022modelsoups}, both trading a
little in-distribution accuracy for out-of-distribution robustness at no extra inference cost. The
theoretical companion most relevant to us is \citet{kumar2022calibrated}, who show that \emph{calibrated
ensembles} --- combining models in \emph{output} space after calibration --- can mitigate the
accuracy--robustness tradeoff under shift; this is the justification behind our composition rule.

Act~II (\Cref{eq:comp}) is deliberately the output-space cousin of these methods: we average the
base's and the adapter's \emph{calibrated scores} with fixed equal weights, retraining nothing. The
tradeoff versus WiSE-FT/soups is explicit --- output-space composition needs only comparable scores,
not interpolable weights, so it is portable across checkpoints that could never be weight-averaged, but
it pays for two inference passes. We position composition as a \emph{transfer-recovery remedy for guard
specialization} specifically (represented cost small, transfer recovered relative to SFT and landing
above base; \Cref{tab:pilot-summary,tab:pilot-models}), and we are explicit about what the pilot does
\emph{not} include --- an actual WiSE-FT weight-space rescoring is
left as a control in the roadmap, and the logit-average variant is reported only as a non-promotable
ablation. The equal-cost SFT+SFT control \emph{is} run (\Cref{tab:sftsft}): base+SFT beats SFT+SFT on every
checkpoint, so the transfer recovery is attributable to keeping the base rather than to generic two-model
ensembling. To our knowledge this calibrated output-space average has not been evaluated as a targeted fix
for the represented/transfer split in small safety guards.

\subsection{Domain-specialized guarding and regulated verticals}
\label{sec:related-domain}
A recent line moves guarding from generic harm to \emph{domain compliance}. ExpGuard
\citep{expguard2026} provides expert-annotated content moderation in specialized domains including
finance, health, and law --- we adopt it directly as our external, expert-labeled breadth replication
(\Cref{tab:expguard}; complete four-checkpoint base result). FinGuard
\citep{dou2026finguard} detects financial regulatory non-compliance in LLM interactions; MortarBench
\citep{toles2026mortarbench} evaluates mortgage loan-origination \emph{agents}; and
\citet{bowen2024mortgage} measure and mitigate racial bias in LLM mortgage \emph{underwriting}.
Adjacent security-benchmark methodology such as Gate AI \citep{goehausen2026gateai} rounds out the
evaluation-design context, and our mortgage instrument is grounded in public HMDA loan-level data
\citep{hmda2022}.

Each of these captures one facet our four-domain design tries to unify while keeping the honesty stance
explicit. FinGuard is single-domain and, like the general guards, reports absolute detection rather than
a paired base-versus-tuned transfer delta. MortarBench evaluates whether an \emph{agent completes} an
origination task, not whether a \emph{guard detects} a compliance violation, and it does not isolate the
requests that read as safe yet violate policy. \citet{bowen2024mortgage} study bias in the model's own
underwriting \emph{decisions}, whereas we study a guard's \emph{detection} behavior and add a
protected-class \emph{minimal-pair} invariance gate as a fairness signal that ranking alone misses. Our
mortgage benchmark's distinctive construct is the \emph{dual label} $G\times D$ (\Cref{tab:composition},
\Cref{fig:pipeline}), whose load-bearing G0/D1 stratum --- looks safe, is a violation --- is exactly the
case a general-safety score cannot see; the zero-shot baselines (\Cref{tab:baseline}) show these
compliance violations are only moderately ranked even by the strongest base. Two honesty boundaries
separate our instruments from the certified-audit reading a reader might infer: the mortgage labels are
\emph{LLM-judge, policy-card-consistent --- not SME-adjudicated}, and ExpGuard supplies the strictly
stronger expert-labeled tier but as a \emph{single-label} transfer probe, not a dual-label construct. We
therefore pair one domain built in depth (mortgage, dual-label plus fairness gate) with three external
verticals for breadth (finance/health/law via ExpGuard), all scored on the same fixed panel used in
Acts~I--II so the specialization question can be asked identically across domains.

% The component-by-component comparison (tab:closest) and its discussion were moved into the
% Introduction (sec:related-closest): the novelty claim is made there, so the table that bounds
% it has to be visible there too rather than 50 pages later.
\noindent The component-by-component comparison against the six closest contributions, and the
narrowed novelty statement it forces, are in \Cref{sec:related-closest} and \Cref{tab:closest} ---
placed in the Introduction beside the claim they bound.

\begin{takeaway}{The gap this report fills}
Prior work establishes, separately, that guards can be built cheaply from small LLMs, that fine-tuning can
degrade or narrow safety, that guards are miscalibrated, that weight-space averaging aids robustness, and
that regulated domains need their own benchmarks. What is missing --- and what
this report supplies on one fixed four-checkpoint panel --- is a single, reproducible, estimation-only
thread that (i) measures the fine-tune as a \emph{paired, same-checkpoint} represented-versus-transfer
delta, (ii) tests a retraining-free \emph{output-space composition} as a targeted transfer-recovery
remedy, and (iii) carries the same question across \emph{four regulated domains} with a dual-label mortgage
construct and a fairness invariance gate. We add no new model or metric; we add the discipline of asking,
at every scale, what the benchmark contributed to the verdict.
\end{takeaway}
